% Cleo bench — shared template for derivation documents.
% Replace <<TITLE>>, <<QID>>, <<N>> and the body. Keep the preamble as-is unless
% a document genuinely needs another package.
\documentclass[11pt]{article}
\usepackage[a4paper,margin=2.7cm]{geometry}
\usepackage{amsmath,amssymb,amsthm,mathtools}
\usepackage[T1]{fontenc}
\usepackage{lmodern}
\usepackage{microtype}
\usepackage{xcolor}
\usepackage[colorlinks=true,linkcolor=blue!50!black,urlcolor=blue!50!black]{hyperref}
\allowdisplaybreaks

\newtheorem{lemma}{Lemma}
\newtheorem{proposition}{Proposition}
\theoremstyle{remark}
\newtheorem*{remark}{Remark}

\title{Cleo Bench Problem 23\\[1ex]\large Integral $\int_0^\infty\frac{\operatorname{arccot}\left(\sqrt{x}-2\,\sqrt{x+1}\right)}{x+1}\mathrm dx$}
\author{Derivation by Claude (Fable 5), closed-book\thanks{Problem originally
posed on Mathematics Stack Exchange
(\href{https://math.stackexchange.com/q/557439}{question 557439}, CC BY-SA),
famously answered by user Cleo. This derivation was produced independently,
offline, without access to the published answer, as part of the Cleo benchmark.}}
\date{July 2026}

\begin{document}
\maketitle

\section*{Problem}

Evaluate in closed form
\[I=\int_0^\infty\frac{\operatorname{arccot}\!\left(\sqrt{x}-2\,\sqrt{x+1}\right)}{x+1}\,\mathrm dx .\]

\textbf{Convention.} For every $x\ge 0$ one has $\sqrt x-2\sqrt{x+1}<0$; in fact (see Step 0) $\sqrt x-2\sqrt{x+1}\le-\sqrt3$. If $\operatorname{arccot}$ were taken with range $(0,\pi)$, the integrand would tend to $\pi/(x+1)$ as $x\to\infty$ and the integral would diverge. Hence the problem must be (and is) read with the \textbf{odd convention}
\[\operatorname{arccot}(y)=\arctan\frac1y\qquad(y\neq0),\]
i.e. the range $(-\tfrac\pi2,0)\cup(0,\tfrac\pi2]$ --- the convention of Mathematica's \texttt{ArcCot} and mpmath's \texttt{acot}. With this convention the integrand is continuous, negative and $O(x^{-3/2})$, so $I$ converges. Everything below uses this convention.

\section*{Result}

\[\boxed{\;\int_0^\infty\frac{\operatorname{arccot}\!\left(\sqrt{x}-2\sqrt{x+1}\right)}{x+1}\,dx
=\frac{\pi}{2}\ln\frac34-\operatorname{Ti}_2\!\left(\frac43\right)
=-\,\pi\ln\frac43-\operatorname{Ti}_2\!\left(\frac34\right)\;}\]

where $\operatorname{Ti}_2(z)=\displaystyle\int_0^z\frac{\arctan t}{t}\,dt=\sum_{k\ge0}\frac{(-1)^k z^{2k+1}}{(2k+1)^2}=\Im\operatorname{Li}_2(iz)$ is Lewin's inverse tangent integral. Numerically
\[I=-1.61434999002764854137373679018430737323694935699855\ldots\]

Equivalent dilogarithmic form: $\;I=\dfrac{\pi}{2}\ln\dfrac34-\Im\operatorname{Li}_2\!\left(\dfrac{4i}{3}\right).$

\section*{Derivation}

\subsection*{Step 0. Convergence and range of the integrand}

Let $y(x)=\sqrt x-2\sqrt{x+1}$. Then $-y(x)=2\sqrt{x+1}-\sqrt x$ attains its minimum where $\frac{d}{dx}(2\sqrt{x+1}-\sqrt x)=\frac{1}{\sqrt{x+1}}-\frac{1}{2\sqrt x}=0$, i.e. $4x=x+1$, $x=\frac13$, with value $2\sqrt{4/3}-\sqrt{1/3}=\sqrt3$. Hence $y(x)\le-\sqrt3$ for all $x\ge0$, so
\[\operatorname{arccot}\big(y(x)\big)=\arctan\frac{1}{y(x)}\in\Big[-\frac\pi6,\,0\Big)\]
is continuous on $[0,\infty)$. As $x\to\infty$, $y(x)=-\sqrt x\,(2\sqrt{1+1/x}-1)\sim-\sqrt x$, so the integrand is $\sim -x^{-3/2}$ and the integral converges absolutely. (With the $(0,\pi)$-range convention the integrand would be $\pi/(x+1)+O(x^{-3/2})$ --- divergent --- which forces the convention stated above.)

\subsection*{Step 1. Trigonometric substitution}

Put $x=\tan^2\theta$, $\theta\in(0,\tfrac\pi2)$. Then $\sqrt{x+1}=\sec\theta$, $x+1=\sec^2\theta$, $dx=2\tan\theta\sec^2\theta\,d\theta$, so $\dfrac{dx}{x+1}=2\tan\theta\,d\theta$, and
\[y=\tan\theta-2\sec\theta=\frac{\sin\theta-2}{\cos\theta}<0 .\]
Therefore, with the odd convention,
\[\operatorname{arccot}(y)=\arctan\frac{\cos\theta}{\sin\theta-2}
=-\arctan\frac{\cos\theta}{2-\sin\theta}=:\eta(\theta),\]
which is smooth on $[0,\tfrac\pi2]$, with $\eta(0)=-\arctan\tfrac12$ and $\eta(\tfrac\pi2)=0$. Hence
\begin{equation}
I=2\int_0^{\pi/2}\eta(\theta)\,\tan\theta\,d\theta. \tag{1}
\end{equation}

\subsection*{Step 2. Differentiate $\eta$ and integrate by parts}

By the quotient rule, with $u(\theta)=\dfrac{\cos\theta}{2-\sin\theta}$:
\[u'=\frac{-\sin\theta\,(2-\sin\theta)+\cos^2\theta}{(2-\sin\theta)^2}
=\frac{1-2\sin\theta}{(2-\sin\theta)^2},\qquad
1+u^2=\frac{(2-\sin\theta)^2+\cos^2\theta}{(2-\sin\theta)^2}=\frac{5-4\sin\theta}{(2-\sin\theta)^2},\]
so
\begin{equation}
\eta'(\theta)=-\frac{u'}{1+u^2}=\frac{2\sin\theta-1}{5-4\sin\theta}. \tag{2}
\end{equation}

Integrate (1) by parts with $\int\tan\theta\,d\theta=-\ln\cos\theta$:
\[I=2\Big[-\eta(\theta)\ln\cos\theta\Big]_0^{\pi/2}+2\int_0^{\pi/2}\eta'(\theta)\ln\cos\theta\,d\theta .\]
Both boundary terms vanish: at $\theta=0$ because $\ln\cos0=0$, and at $\theta\to\tfrac\pi2$ because
\[|\eta(\theta)\ln\cos\theta|\le\frac{\cos\theta}{2-\sin\theta}\,|\ln\cos\theta|\le\cos\theta\,|\ln\cos\theta|\longrightarrow0\]
(using $|\arctan v|\le|v|$). Hence, by (2),
\begin{equation}
I=2\int_0^{\pi/2}\frac{(2\sin\theta-1)\ln\cos\theta}{5-4\sin\theta}\,d\theta. \tag{3}
\end{equation}

\subsection*{Step 3. Split off the elementary part}

Since $2\sin\theta-1=-\tfrac12(5-4\sin\theta)+\tfrac32$,
\begin{equation}
I=-\int_0^{\pi/2}\ln\cos\theta\,d\theta+3\int_0^{\pi/2}\frac{\ln\cos\theta}{5-4\sin\theta}\,d\theta
=\frac\pi2\ln2+3W,\qquad W:=\int_0^{\pi/2}\frac{\ln\cos\theta}{5-4\sin\theta}\,d\theta, \tag{4}
\end{equation}
where we used the classical Euler integral $\int_0^{\pi/2}\ln\cos\theta\,d\theta=-\frac\pi2\ln2$. (Proof: $J:=\int_0^{\pi/2}\ln\sin=\int_0^{\pi/2}\ln\cos$ by $\theta\mapsto\frac\pi2-\theta$; then $2J=\int_0^{\pi/2}\ln\frac{\sin2\theta}{2}\,d\theta=\frac12\int_0^{\pi}\ln\sin t\,dt-\frac\pi2\ln2=J-\frac\pi2\ln2$.)

\subsection*{Step 4. Evaluate $W$: rationalize the kernel}

Multiply numerator and denominator by $5+4\sin\theta$ and use $25-16\sin^2\theta=9+16\cos^2\theta$:
\begin{equation}
\frac{1}{5-4\sin\theta}=\frac{5+4\sin\theta}{9+16\cos^2\theta}
\qquad\Longrightarrow\qquad
W=5\,W_1+4\,W_2, \tag{5}
\end{equation}
\[W_1:=\int_0^{\pi/2}\frac{\ln\cos\theta}{9+16\cos^2\theta}\,d\theta,\qquad
W_2:=\int_0^{\pi/2}\frac{\sin\theta\,\ln\cos\theta}{9+16\cos^2\theta}\,d\theta .\]

\textbf{Lemma A.} For $p\ge0$: $\displaystyle\int_0^\infty\frac{\ln(1+p^2v^2)}{1+v^2}\,dv=\pi\ln(1+p).$

\textit{Proof.} Call it $J(p)$; $J(0)=0$. For $p$ in a compact subset of $(0,\infty)$ the differentiated integrand $\frac{2pv^2}{(1+p^2v^2)(1+v^2)}$ is dominated by $\frac{C}{1+v^2}\cdot$const uniformly, so differentiation under the integral sign is legitimate (dominated convergence), and by the partial fraction $\frac{v^2}{(1+p^2v^2)(1+v^2)}=\frac{1}{p^2-1}\big[\frac{1}{1+v^2}-\frac{1}{1+p^2v^2}\big]$ (for $p\neq1$; $p=1$ removable by continuity),
\[J'(p)=\frac{2p}{p^2-1}\Big[\frac\pi2-\frac{\pi}{2p}\Big]=\frac{\pi}{p+1}.\]
Integrating from $0$: $J(p)=\pi\ln(1+p)$. $\blacksquare$

\textbf{Evaluation of $W_1$.} Write $9+16\cos^2\theta=9\sin^2\theta+25\cos^2\theta$ and substitute $v=\tan\theta$ ($\cos^2\theta=\frac{1}{1+v^2}$, $\ln\cos\theta=-\frac12\ln(1+v^2)$, $d\theta=\frac{dv}{1+v^2}$):
\[W_1=\int_0^{\pi/2}\frac{\ln\cos\theta}{\cos^2\theta\,(9\tan^2\theta+25)}\,\frac{d\theta\,\cos^2\theta}{1}
=-\frac12\int_0^\infty\frac{\ln(1+v^2)}{9v^2+25}\,dv .\]
Now $v=\frac53 w$ gives $\int_0^\infty\frac{\ln(1+v^2)}{9v^2+25}dv=\frac{1}{15}\int_0^\infty\frac{\ln\!\big(1+\frac{25}{9}w^2\big)}{1+w^2}dw=\frac{\pi}{15}\ln\frac83$ by Lemma A with $p=\frac53$. Hence
\begin{equation}
W_1=-\frac{\pi}{30}\ln\frac83. \tag{6}
\end{equation}

\textbf{Lemma B.} For $y>0$: $\displaystyle\int_0^{y}\frac{\ln v}{1+v^2}\,dv=\arctan y\,\ln y-\operatorname{Ti}_2(y).$

\textit{Proof.} Integrate by parts, $d(\arctan v)=\frac{dv}{1+v^2}$: the boundary term $\arctan v\,\ln v\to0$ as $v\to0^+$ (since $\arctan v\sim v$), and $\int_0^y\frac{\arctan v}{v}dv=\operatorname{Ti}_2(y)$ by definition. $\blacksquare$

\textbf{Evaluation of $W_2$.} Substitute $t=\cos\theta$:
\[W_2=\int_0^1\frac{\ln t}{9+16t^2}\,dt
\overset{t=\frac34 v}{=}\frac{1}{12}\int_0^{4/3}\frac{\ln\frac34+\ln v}{1+v^2}\,dv
=\frac1{12}\Big[\ln\tfrac34\,\arctan\tfrac43+\arctan\tfrac43\,\ln\tfrac43-\operatorname{Ti}_2\big(\tfrac43\big)\Big],\]
by Lemma B with $y=\frac43$. The two logarithmic terms cancel ($\ln\frac34+\ln\frac43=0$), so
\begin{equation}
W_2=-\frac{1}{12}\operatorname{Ti}_2\!\left(\frac43\right). \tag{7}
\end{equation}

\subsection*{Step 5. Assemble}

From (5), (6), (7):
\[W=-\frac{\pi}{6}\ln\frac83-\frac13\operatorname{Ti}_2\!\left(\frac43\right),\]
and from (4):
\begin{equation}
I=\frac{\pi}{2}\ln2-\frac{\pi}{2}\ln\frac83-\operatorname{Ti}_2\!\left(\frac43\right)
=\frac{\pi}{2}\ln\frac{3}{4}-\operatorname{Ti}_2\!\left(\frac43\right). \tag{8}
\end{equation}

\subsection*{Step 6. Equivalent forms}

\textbf{Lemma C (inversion).} For $y>0$: $\operatorname{Ti}_2(y)-\operatorname{Ti}_2(1/y)=\frac\pi2\ln y$.

\textit{Proof.} Both sides vanish at $y=1$ and have equal derivatives: $\frac{d}{dy}\big[\operatorname{Ti}_2(y)-\operatorname{Ti}_2(1/y)\big]=\frac{\arctan y}{y}+\frac{\arctan(1/y)}{y}=\frac{\pi}{2y}$, using $\arctan y+\arctan\frac1y=\frac\pi2$ for $y>0$. $\blacksquare$

With $y=\frac43$, (8) becomes
\[I=\frac\pi2\ln\frac34-\frac\pi2\ln\frac43-\operatorname{Ti}_2\!\left(\frac34\right)
=-\pi\ln\frac43-\operatorname{Ti}_2\!\left(\frac34\right).\]
Also $\operatorname{Ti}_2(y)=\Im\operatorname{Li}_2(iy)$ for $y>0$ (the Maclaurin identity $\operatorname{Li}_2(iy)=\frac14\operatorname{Li}_2(-y^2)+i\operatorname{Ti}_2(y)$ for $y\le1$ extends to all $y>0$ by analytic continuation along the imaginary axis, which never meets the cut $[1,\infty)$ of $\operatorname{Li}_2$), giving
\[I=\frac{\pi}{2}\ln\frac34-\Im\operatorname{Li}_2\!\left(\frac{4i}{3}\right)=-\pi\ln\frac43-\Im\operatorname{Li}_2\!\left(\frac{3i}{4}\right).\]

\textit{Remark on the shape of the answer.} Since $\arctan\frac43=2\arctan\frac12$ is not a rational multiple of $\pi$, the standard reduction $\operatorname{Ti}_2(\tan\phi)=\phi\ln\tan\phi+\frac12\operatorname{Cl}_2(2\phi)+\frac12\operatorname{Cl}_2(\pi-2\phi)$ produces no further elementary simplification; $\operatorname{Ti}_2\!\big(\frac43\big)$ (equivalently $\operatorname{Ti}_2\!\big(\frac34\big)$) is the irreducible transcendental constant of this problem, exactly as Catalan's constant $G=\operatorname{Ti}_2(1)$ would be. See the Notes for a PSLQ search confirming no small-height relation with the usual constant basis.

\section*{Numerical verification}

All computations with mpmath.

\begin{enumerate}
\item \textbf{Direct evaluation of the integral as posed} (odd-convention \texttt{acot}, i.e. \texttt{mp.acot}), at \texttt{mp.dps = 80}, splitting $[0,\infty)$ at $x=\frac13,1,4,20,100$ and mapping the tail $x>100$ by $x=1/t^2$:
\[I_{\text{quad}}=-1.6143499900276485413737367901843073732369493569985507490797150\ldots\]
\item \textbf{Independent direct evaluation} in the form (1) (pure change of variable, before any integration by parts), $2\int_0^{\pi/2}\eta\tan\theta\,d\theta$, at \texttt{mp.dps = 60}: agrees with the closed form to $\sim6\times10^{-61}$ (all 60 digits).
\item \textbf{Closed form} $\frac\pi2\ln\frac34-\operatorname{Ti}_2(\frac43)$ with $\operatorname{Ti}_2(\frac43)=\int_0^{4/3}\frac{\arctan t}{t}dt$ by quadrature, and independently as $\Im\operatorname{Li}_2(4i/3)$ and $-\pi\ln\frac43-\Im\operatorname{Li}_2(0.75i)$ via \texttt{mp.polylog}; all four forms agree among themselves to full working precision ($\le3\times10^{-61}$ at 60 dps).
\item \textbf{Agreement:} $|I_{\text{quad}}-I_{\text{closed}}|\approx2.1\times10^{-81}$ at 80 digits of working precision --- \textbf{80 matching significant digits}.
\item Every intermediate identity was checked numerically at 60 digits: the derivative formula (2), equation (3), the value of $W$, $W_1$ vs $-\frac{\pi}{30}\ln\frac83$, $W_2$ vs $-\frac1{12}\operatorname{Ti}_2(\frac43)$, Lemma A at $p=\frac53$, and Lemma C at $y=\frac43$. All residuals were $<10^{-58}$.
\end{enumerate}

Reference values (60 significant digits):
\[I=-1.61434999002764854137373679018430737323694935699855074907972\ldots\]
\[\operatorname{Ti}_2\!\left(\tfrac43\right)=1.16246004733564774158995818385573624847446067214211087965611\ldots\]

\section*{Notes}

\begin{itemize}
\item \textbf{Convention caveat.} The result assumes $\operatorname{arccot} y=\arctan\frac1y$ (odd convention, range $(-\frac\pi2,0)\cup(0,\frac\pi2]$). This is forced: with the continuous $(0,\pi)$ convention the integral diverges to $+\infty$. This is the standard reading for such forum problems (Mathematica's \texttt{ArcCot}).
\item \textbf{Irreducibility of $\operatorname{Ti}_2(4/3)$.} A PSLQ search at 60 digits for integer relations (coefficients up to $10^6$--$10^8$) between $\operatorname{Ti}_2(\frac34)$ and the basis $\{G,\ \pi\ln2,\ \pi\ln3,\ \pi\ln5,\ \arctan\!\frac12\ln p,\ \arctan\!\frac13\ln p\ (p=2,3,5),\ \pi^2\}$ found only the trivial relation coming from $\arctan\frac12+\arctan\frac13=\frac\pi4$. So the answer very likely cannot be reduced to $G$, $\pi$, and logarithms; $\operatorname{Ti}_2(\frac43)$ is the natural final constant. (Known closed-form points of $\operatorname{Ti}_2$ are $\tan\phi$ with $\phi$ a special rational multiple of $\pi$: $1,\ 2\pm\sqrt3$, etc.; $\frac43$ is not of this type.)
\item \textbf{A curious empirical relation (not used above).} PSLQ does find, and I verified to 150 digits,
\[\operatorname{Ti}_2\!\left(\tfrac34\right)+4\operatorname{Ti}_2\!\left(\tfrac12\right)+2\operatorname{Ti}_2\!\left(\tfrac13\right)=6G-\pi\ln2,\]
a Lewin-type multi-argument relation (consistent with $2\arctan\frac13=\arctan\frac34$ and $\arctan\frac12+\arctan\frac13=\frac\pi4$). I did not prove it and the main result does not depend on it; it shows the answer can also be written as $I=-\pi\ln\frac43-6G+4\operatorname{Ti}_2(\frac12)+2\operatorname{Ti}_2(\frac13)+\pi\ln2$, which is no simpler.
\item \textbf{Rigor status.} The derivation is complete and elementary: one change of variables, one integration by parts with vanishing boundary terms (justified), a rational kernel split, Lemma A (differentiation under the integral, justified by dominated convergence), Lemma B (integration by parts), Lemma C (inversion). No steps are heuristic. Combined with 80-digit numerical agreement, I consider the result established.
\end{itemize}

\end{document}
